\chapter{Fourier Duality, Momentum, and Free-Particle Motion}

\section{From Foundations to Examples}

Entanglement belongs at the heart of quantum mechanics, but the foundations are now sufficiently exposed that we can begin using them. The course has so far relied on deliberately small systems: a spin, a rotatable apparatus, a finite list of outcomes, and repeated measurements from which probabilities are inferred. Those experiments are logically simpler than the historical path by which quantum mechanics was discovered, but they contain the minimum operational ingredients.

\begin{classroomqa}
\audiencequestion{What plays the role of Newton's small set of empirical starting points in quantum mechanics?}
\lecturerresponse{A minimal model is the rotatable spin apparatus. Prepare a spin, choose an axis, press the measurement control, record \(+1\) or \(-1\), and repeat the experiment until the outcome frequencies are known. The vector space, observables, amplitudes, and Born rule organize those statistics. Historically the theory arose from a much richer collection of atomic phenomena, but logically this two-outcome experiment is enough to expose its structure.}
\end{classroomqa}

\begin{classroomqa}
\audiencequestion{Must the same statistical procedure be carried out for each physical quantity, such as spin or momentum?}
\lecturerresponse{Yes. An observable is tied to a class of experimental arrangements and possible outcomes. The formalism predicts the distribution once the state and observable are specified, but the physical meaning of the observable is anchored by how it is measured.}
\end{classroomqa}

Before leaving entanglement, one question from the previous discussion deserves a complete answer. Consider two electrons with singlet and triplet sectors and include the radiation field as another quantum subsystem. The radiation field has a vacuum state \(\ket{0_\gamma}\) and a one-photon state \(\ket{1_\gamma}\). Particle number itself may be superposed; there is no rule restricting superposition to different states of a fixed number of particles.

For the idealized transition discussed in class, an initially unentangled electron state may be written as a superposition of singlet and triplet components,
\begin{equation}
\ket{\uparrow\downarrow}
=
\frac{1}{\sqrt{2}}
\left(\ket{S}+\ket{T_0}\right).
\end{equation}
With no photon initially present, the state is
\begin{equation}
\ket{\Psi(0)}
=
\frac{1}{\sqrt{2}}
\left(\ket{S}+\ket{T_0}\right)
\otimes\ket{0_\gamma}.
\label{eq:l9-initial-radiative-state}
\end{equation}
The singlet component remains a singlet without emission, while the excited triplet component can relax to the singlet and emit a photon:
\begin{align}
\ket{S}\ket{0_\gamma}
&\longrightarrow
\ket{S}\ket{0_\gamma}, \\
\ket{T_0}\ket{0_\gamma}
&\longrightarrow
\ket{S}\ket{1_\gamma}.
\end{align}
Linearity then gives
\begin{equation}
\ket{\Psi_{\mathrm{final}}}
=
\ket{S}\otimes
\frac{\ket{0_\gamma}+\ket{1_\gamma}}{\sqrt{2}}.
\label{eq:l9-final-radiative-state}
\end{equation}

\begin{figure}[htbp]
\centering
\includegraphics[width=0.88\textwidth]{lecture_09_figure_01.png}
\caption{The electron state is decomposed into singlet and triplet components and tensored with the photon vacuum before radiative relaxation.}
\lectureframe{00:05:10}
\end{figure}

The final expression is a product between the electron pair and the radiation field. The two electrons are entangled with each other because \(\ket{S}\) is entangled, while the radiation field is left in a superposition of zero and one photon. This example is subtle because interaction has redistributed entanglement: the initial electron product state evolves into an entangled electron pair, yet the final idealized electron-radiation state factorizes.

\begin{classroomqa}
\audiencequestion{Could any incoming photon reverse the transition and disentangle the electrons?}
\lecturerresponse{No. The time-reversed process requires the appropriate photon state, mode, phase, energy, and interaction history. Reflecting the emitted photon back was shorthand for arranging the inverse evolution of the entire closed system, not for claiming that an arbitrary photon will undo the entanglement.}
\end{classroomqa}

\begin{classroomqa}
\audiencequestion{What causes the triplet to emit if no external magnetic field has been introduced?}
\lecturerresponse{The electrons interact with one another and with the quantized electromagnetic field. Their spin-dependent interaction makes one spin configuration energetically favorable, and the excess energy can leave as radiation. No external measuring apparatus is needed; this is ordinary unitary dynamics of the combined electrons and field.}
\end{classroomqa}

With that remaining point settled, we now forget spin and study the simplest motion problem: one particle on a one-dimensional line.

\section{Plane Waves and the Fourier Pair}

Before returning to physics, we need a mathematical interlude. Begin with
\begin{equation}
e^{ipx}=\cos(px)+i\sin(px).
\end{equation}
At this stage \(p\) is simply a real number that controls the rate of oscillation. The expression has no time dependence and therefore does not yet describe a moving wave. It is called a plane wave because, as a function of \(x\), it has the familiar sinusoidal form.

A broad class of sufficiently regular functions can be decomposed into such modes. We restrict attention to functions that behave well enough for the integrals below: they are not allowed to blow up at spatial infinity, and later physical wave functions will be square-integrable. With a symmetric normalization convention, the Fourier pair is
\begin{align}
\psi(x)
&=
\int_{-\infty}^{\infty}
\frac{dp}{\sqrt{2\pi}}\,
\widetilde{\psi}(p)e^{ipx},
\label{eq:l9-fourier-forward}\\
\widetilde{\psi}(p)
&=
\int_{-\infty}^{\infty}
\frac{dx}{\sqrt{2\pi}}\,
\psi(x)e^{-ipx}.
\label{eq:l9-fourier-inverse}
\end{align}
The coefficient \(\widetilde{\psi}(p)\) is the complex amplitude of the mode with Fourier label \(p\). The factors of \(1/\sqrt{2\pi}\) are conventional; placing one in each direction makes the two formulas look reciprocal. The relative minus sign in the exponent is essential.

\begin{figure}[htbp]
\centering
\includegraphics[width=0.88\textwidth]{lecture_09_figure_02.png}
\caption{The symmetric Fourier pair: a function is decomposed into plane waves, and the mode amplitudes are recovered by the inverse transform.}
\lectureframe{00:14:35}
\end{figure}

\begin{classroomqa}
\audiencequestion{Are \(\psi(x)\) and \(\widetilde{\psi}(p)\) intended to be complex?}
\lecturerresponse{In general, yes. A special function may happen to be real, in which case its Fourier coefficients obey an additional conjugation relation, but quantum amplitudes are naturally complex and no reality restriction is imposed here.}
\end{classroomqa}

\begin{classroomqa}
\audiencequestion{Why use the symmetric factors \(1/\sqrt{2\pi}\) rather than put all of \(2\pi\) in one transform?}
\lecturerresponse{Either convention is consistent. The symmetric choice makes the forward and inverse transforms differ only by the sign in the exponential. Physical predictions do not depend on the convention provided every related normalization is changed consistently.}
\end{classroomqa}

The reciprocal relation between \(x\) and \(p\) is already visible. One function is resolved by integrating over the other variable. For the moment this is Fourier analysis. Soon the same pair will become position and momentum.

\section{The Dirac Delta as an Operation}

The intuitive delta function is a spike of vanishing width, diverging height, and unit area. That picture is useful, but the operational definition is more reliable. For every suitable test function \(f\),
\begin{equation}
\int_{-\infty}^{\infty}
dx\,\delta(x-x')f(x)
=f(x').
\label{eq:l9-delta-sampling}
\end{equation}
Three familiar rules express the same idea:
\begin{align}
f(x)\delta(x-x')
&=
f(x')\delta(x-x'),
\label{eq:l9-delta-multiplication}\\
\int_{-\infty}^{\infty}dx\,\delta(x-x')
&=1,
\label{eq:l9-delta-area}\\
\int_{-\infty}^{\infty}dx\,\delta(x-x')f(x)
&=f(x').
\end{align}
The first rule holds because the delta has support only where \(x=x'\). In the second rule, translating the spike changes its location but not its area. Combining those two rules gives the third.

\begin{figure}[htbp]
\centering
\includegraphics[width=0.88\textwidth]{lecture_09_figure_03.png}
\caption{The sampling property that characterizes the Dirac delta by its action on every suitable function.}
\lectureframe{00:19:45}
\end{figure}

\begin{classroomqa}
\audiencequestion{Why may \(f(x')\) be pulled outside the \(x\)-integral?}
\lecturerresponse{The integration variable is \(x\). The label \(x'\) is fixed during that integration, so \(f(x')\) is a constant with respect to \(x\). The remaining delta integral has unit area.}
\end{classroomqa}

The delta is not an ordinary pointwise function. In particular, trying to square it leads to an object more singular than the original delta. Its safe meaning is distributional: Equation~\eqref{eq:l9-delta-sampling} specifies what it does under an integral.

\subsection{The Fourier representation of the delta}

Insert Equation~\eqref{eq:l9-fourier-inverse} into Equation~\eqref{eq:l9-fourier-forward}, renaming the argument on the left \(x'\):
\begin{align}
\psi(x')
&=
\int\frac{dp}{\sqrt{2\pi}}\,
e^{ipx'}
\left[
\int\frac{dx}{\sqrt{2\pi}}\,
\psi(x)e^{-ipx}
\right] \notag\\
&=
\int dx
\left[
\int\frac{dp}{2\pi}
e^{ip(x'-x)}
\right]\psi(x).
\label{eq:l9-fourier-composition}
\end{align}
The expression in square brackets must reproduce \(\psi(x')\) for every suitable \(\psi\). Comparing with the defining action of the delta gives
\begin{equation}
\boxed{
\int_{-\infty}^{\infty}
\frac{dp}{2\pi}\,
e^{ip(x'-x)}
=\delta(x'-x)
}.
\label{eq:l9-delta-fourier-x}
\end{equation}

\begin{figure}[htbp]
\centering
\includegraphics[width=0.88\textwidth]{lecture_09_figure_04.png}
\caption{The Fourier kernel is isolated inside the composition of the two transforms and identified by its sampling action as a Dirac delta.}
\lectureframe{00:26:23}
\end{figure}

Interchanging the names of the reciprocal variables yields the companion identity
\begin{equation}
\boxed{
\int_{-\infty}^{\infty}
\frac{dx}{2\pi}\,
e^{ix(p-p')}
=\delta(p-p')
}.
\label{eq:l9-delta-fourier-p}
\end{equation}
The delta function is thus a continuous superposition of plane waves.

\begin{classroomqa}
\audiencequestion{The exponent was written once as \(x'-x\) and elsewhere as \(x-x'\). Does the sign change alter the result?}
\lecturerresponse{No. The delta is even: \(\delta(u)=\delta(-u)\). Reversing the sign of the difference therefore leaves the final distribution unchanged. The intermediate exponential changes to its complex conjugate, but the integrated kernel is the same delta.}
\end{classroomqa}

These manipulations suppress a convergence qualification. The oscillatory integrals are distributions, not absolutely convergent ordinary integrals, and changing the order of integration requires care. A finite periodic interval supplies a clean regulator: first use discrete Fourier modes and Kronecker orthogonality, then let the interval grow without bound. The continuum identities above are the resulting limit.

\section{A Particle on a Line}

Classically, one point in phase space specifies the state:
\begin{equation}
(x,p).
\end{equation}
Quantum mechanics asks a different first question: what orthonormal basis defines the state space? For each real position \(x\), introduce a generalized eigenket \(\ket{x}\). A general state can be expanded as
\begin{equation}
\ket{\Psi}
=
\int_{-\infty}^{\infty}
dx\,\psi(x)\ket{x},
\qquad
\psi(x)=\braket{x}{\Psi}.
\label{eq:l9-position-expansion}
\end{equation}

\begin{figure}[htbp]
\centering
\includegraphics[width=0.88\textwidth]{lecture_09_figure_05.png}
\caption{A general state expanded in the continuous position basis, with the coefficient recovered by projection onto \(\ket{x}\).}
\lectureframe{00:35:30}
\end{figure}

The Born rule assigns the probability density
\begin{equation}
\rho(x)=|\psi(x)|^2.
\end{equation}
The probability in a finite interval is
\begin{equation}
P(a<x<b)=\int_a^b|\psi(x)|^2\,dx,
\end{equation}
and a normalized physical state satisfies
\begin{equation}
\int_{-\infty}^{\infty}|\psi(x)|^2\,dx=1.
\label{eq:l9-position-normalization}
\end{equation}
Continuous orthogonality is expressed by
\begin{equation}
\braket{x'}{x}=\delta(x-x').
\label{eq:l9-position-delta}
\end{equation}
For two arbitrary states, the inner product becomes
\begin{equation}
\braket{\Psi}{\Phi}
=
\int_{-\infty}^{\infty}
dx\,\psi^*(x)\phi(x).
\label{eq:l9-position-inner-product}
\end{equation}

\begin{figure}[htbp]
\centering
\includegraphics[width=0.88\textwidth]{lecture_09_figure_06.png}
\caption{Continuous basis normalization and the wave-function inner product written in the position representation.}
\lectureframe{00:41:20}
\end{figure}

\begin{classroomqa}
\audiencequestion{Why is position alone enough? Would adding momentum or almost any other variable be redundant?}
\lecturerresponse{For one particle in one dimension, the complete function \(\psi(x)\) specifies the state, so a simultaneous sharp value of \(p\) would be too much information. Momentum supplies another complete basis rather than an extra label attached to each \(\ket{x}\). In more dimensions, compatible information can be mixed: for example, \(X\) and \(P_y\) commute and may label a joint basis. The restriction concerns noncommuting pairs, not the mere number of symbols.}
\end{classroomqa}

\begin{classroomqa}
\audiencequestion{Can momentum also supply a complete basis and be expanded through position?}
\lecturerresponse{Yes. The momentum basis is complete, and its coefficients are related to the position coefficients by Fourier transformation. That relation is the central construction of this lecture.}
\end{classroomqa}

\section{Position in the Position Representation}

An abstract operator acts on state vectors. Once a basis is chosen, the same action becomes a concrete operation on wave functions. In the position representation, the position operator \(X\) acts by multiplication:
\begin{equation}
(X\psi)(x)=x\psi(x).
\label{eq:l9-position-operator}
\end{equation}

The state localized at \(x_0\) is represented by the generalized wave function \(\delta(x-x_0)\). Applying \(X\) gives
\begin{align}
X\delta(x-x_0)
&=
x\delta(x-x_0) \notag\\
&=
x_0\delta(x-x_0),
\label{eq:l9-position-eigenfunction}
\end{align}
where Equation~\eqref{eq:l9-delta-multiplication} was used. Thus \(\delta(x-x_0)\) is an eigenfunction of \(X\) with eigenvalue \(x_0\).

\begin{figure}[htbp]
\centering
\includegraphics[width=0.88\textwidth]{lecture_09_figure_07.png}
\caption{Multiplication by \(x\) turns the delta-localized state into \(x_0\) times itself, proving that it is a position eigenfunction.}
\lectureframe{00:45:50}
\end{figure}

Neither \(\ket{x_0}\) nor its delta-function representative is normalized to one in the ordinary square-integrable sense. It is delta-normalized. Physical localized states are finite wave packets, while the exact eigenstates serve as generalized basis vectors.

\section{Momentum Eigenfunctions}

The previous lecture found that differentiation is anti-Hermitian and that multiplication by \(-i\) makes it Hermitian. Restoring units, define
\begin{equation}
P=-i\hbar\frac{d}{dx},
\qquad
(P\psi)(x)
=
-i\hbar\frac{d\psi}{dx}.
\label{eq:l9-momentum-operator}
\end{equation}
This is a legitimate observable on a suitable domain. Its identification with physical momentum will be tested through its spectrum, wavelength, and role in dynamics.

For a momentum eigenvalue \(p\), solve
\begin{equation}
P\psi_p=p\psi_p,
\end{equation}
or
\begin{equation}
-i\hbar\frac{d\psi_p}{dx}=p\psi_p.
\end{equation}
The solution is
\begin{equation}
\psi_p(x)=A e^{ipx/\hbar}.
\label{eq:l9-momentum-eigenfunction}
\end{equation}

\begin{figure}[htbp]
\centering
\includegraphics[width=0.88\textwidth]{lecture_09_figure_08.png}
\caption{The position-space representative of a definite-momentum state is a plane wave \(e^{ipx/\hbar}\).}
\lectureframe{00:51:45}
\end{figure}

The wavelength \(\lambda\) is the displacement that advances the phase by \(2\pi\):
\begin{equation}
\frac{p\lambda}{\hbar}=2\pi,
\qquad
\lambda=\frac{2\pi\hbar}{p}.
\label{eq:l9-de-broglie}
\end{equation}
Greater momentum therefore corresponds to shorter wavelength.

Delta normalization fixes the constant:
\begin{equation}
\psi_p(x)
=
\frac{1}{\sqrt{2\pi\hbar}}e^{ipx/\hbar}.
\label{eq:l9-normalized-momentum-eigenfunction}
\end{equation}
Indeed,
\begin{align}
\braket{p'}{p}
&=
\int_{-\infty}^{\infty}
dx\,
\psi_{p'}^*(x)\psi_p(x) \notag\\
&=
\int_{-\infty}^{\infty}
\frac{dx}{2\pi\hbar}
e^{i(p-p')x/\hbar} \notag\\
&=
\delta(p-p').
\label{eq:l9-momentum-normalization}
\end{align}
After choosing units with \(\hbar=1\), these formulas become
\begin{equation}
\psi_p(x)=\frac{1}{\sqrt{2\pi}}e^{ipx},
\qquad
\int\frac{dx}{2\pi}e^{i(p-p')x}
=\delta(p-p').
\end{equation}

\begin{classroomqa}
\audiencequestion{How can the oscillatory integral be declared zero for \(p\neq p'\) when it does not converge in the ordinary sense?}
\lecturerresponse{The lecture takes a controlled shortcut. Put the system on a finite periodic interval. The allowed momenta then form a discrete set, the exponentials are ordinarily orthogonal, and every normalization is finite. Only after that calculation should the interval be sent to infinity, turning the Kronecker delta into a Dirac delta.}
\end{classroomqa}

\begin{classroomqa}
\audiencequestion{Does the same convergence issue affect the Fourier-transform derivation of the delta function?}
\lecturerresponse{Yes. Interchanging the integrations there was formal for the same reason. The finite-box derivation or the language of distributions supplies the omitted rigor. The shortcut is acceptable because the regulated calculation has the stated continuum limit.}
\end{classroomqa}

\section{The Momentum Representation}

Two objects must not be confused. The function \(\psi_p(x)\) in Equation~\eqref{eq:l9-normalized-momentum-eigenfunction} is the position-space representative of the particular eigenstate \(\ket{p}\). By contrast, the momentum-space wave function of an arbitrary state \(\ket{\Psi}\) is
\begin{equation}
\widetilde{\psi}(p)=\braket{p}{\Psi}.
\label{eq:l9-momentum-wavefunction}
\end{equation}
Insert the position-space wave functions of \(\bra{p}\) and \(\ket{\Psi}\):
\begin{align}
\widetilde{\psi}(p)
&=
\int_{-\infty}^{\infty}
dx\,
\braket{p}{x}\braket{x}{\Psi} \notag\\
&=
\int_{-\infty}^{\infty}
\frac{dx}{\sqrt{2\pi\hbar}}\,
e^{-ipx/\hbar}\psi(x).
\label{eq:l9-position-to-momentum}
\end{align}
This is precisely the Fourier transform. Conversely,
\begin{equation}
\psi(x)
=
\int_{-\infty}^{\infty}
\frac{dp}{\sqrt{2\pi\hbar}}\,
\widetilde{\psi}(p)e^{ipx/\hbar}.
\label{eq:l9-momentum-to-position}
\end{equation}

\begin{figure}[htbp]
\centering
\includegraphics[width=0.88\textwidth]{lecture_09_figure_09.png}
\caption{Projection onto \(\ket{p}\) turns the position-space wave function into its Fourier transform \(\widetilde{\psi}(p)\).}
\lectureframe{01:05:50}
\end{figure}

The probability density for measuring momentum \(p\) is
\begin{equation}
\widetilde{\rho}(p)
=
|\widetilde{\psi}(p)|^2,
\qquad
\int_{-\infty}^{\infty}
|\widetilde{\psi}(p)|^2\,dp=1.
\end{equation}
Thus one position-space wave function contains enough information to predict both position and momentum measurements. To obtain momentum probabilities, Fourier transform \(\psi(x)\) and square the modulus of the result.

In the momentum representation, the roles reverse:
\begin{equation}
(P\widetilde{\psi})(p)
=
p\widetilde{\psi}(p),
\qquad
(X\widetilde{\psi})(p)
=
i\hbar\frac{d\widetilde{\psi}}{dp}.
\label{eq:l9-operators-in-p}
\end{equation}
The sign difference from \(P=-i\hbar\,d/dx\) follows from the opposite signs in the Fourier pair.

\begin{classroomqa}
\audiencequestion{Is time the operator conjugate to energy in the same way that position is conjugate to momentum?}
\lecturerresponse{Time is usually treated as the external parameter labeling state evolution, not as a universal self-adjoint operator on the system's Hilbert space. A physical clock does have observable readings, and improving the time resolution of a clock has energy consequences, but the time-energy relation requires more care than the position-momentum pair and is not developed here.}
\end{classroomqa}

\begin{classroomqa}
\audiencequestion{What other conjugate pairs occur?}
\lecturerresponse{Angle and the corresponding angular momentum provide an important example, with global qualifications because angle is periodic. In field theory, a field and its canonical momentum are conjugate. Temperature and pressure are thermodynamic variables rather than a canonical quantum pair. The general principle is supplied by canonical variables and their commutators, not by any arbitrary pair of measurable quantities.}
\end{classroomqa}

\begin{classroomqa}
\audiencequestion{In the discrete theory operators acted on kets and wave functions merely listed coefficients. Has that changed in the continuous theory?}
\lecturerresponse{No. A wave function is the coordinate column of an abstract ket in a chosen basis. A matrix acting on a finite column and a differential or integral operator acting on a function are two concrete representations of the same abstract operation. The continuous label changes the notation, not the logic.}
\end{classroomqa}

\subsection{Uncertainty as a Fourier fact}

A definite-momentum state is completely delocalized:
\begin{equation}
|\psi_p(x)|^2=\text{constant}.
\end{equation}
Its momentum is sharp, but its position is spread uniformly across the line. More generally, Fourier pairs cannot both be narrow. Concentrating \(\widetilde{\psi}(p)\) into a small range of momenta forces \(\psi(x)\) to extend over a large region, and localizing \(\psi(x)\) requires a broad range of momenta.

The lecture does not yet derive the quantitative inequality
\begin{equation}
\Delta X\,\Delta P\geq\frac{\hbar}{2}.
\end{equation}
Its present point is structural: the uncertainty principle is already encoded in the reciprocal geometry of Fourier transforms. This is the continuous analogue of changing between incompatible spin bases.

\section{The Hamiltonian Generates Time Evolution}

The Hamiltonian \(H\) is Hermitian, represents energy, and generates unitary time evolution. The time-dependent Schr\"odinger equation is
\begin{equation}
i\hbar\frac{d}{dt}\ket{\Psi(t)}
=
H\ket{\Psi(t)}.
\label{eq:l9-abstract-schrodinger}
\end{equation}
In the position representation,
\begin{equation}
i\hbar\frac{\partial\psi(x,t)}{\partial t}
=
H\psi(x,t),
\label{eq:l9-position-schrodinger}
\end{equation}
where the right side means the chosen position-space representation of the Hamiltonian acting on the function.

\subsection{A deliberately simple linear Hamiltonian}

Take
\begin{equation}
H=cP.
\label{eq:l9-linear-hamiltonian}
\end{equation}
This is not the ordinary nonrelativistic free-particle Hamiltonian. It is chosen because every step can be seen directly.

\begin{figure}[htbp]
\centering
\includegraphics[width=0.88\textwidth]{lecture_09_figure_10.png}
\caption{The Schr\"odinger equation and the linear Hamiltonian \(H=cP\), chosen as the simplest example of particle evolution.}
\lectureframe{01:17:45}
\end{figure}

Using \(P=-i\hbar\,\partial_x\),
\begin{align}
i\hbar\frac{\partial\psi}{\partial t}
&=
c\left(-i\hbar\frac{\partial}{\partial x}\right)\psi, \notag\\
\frac{\partial\psi}{\partial t}
&=
-c\frac{\partial\psi}{\partial x},
\end{align}
or
\begin{equation}
\frac{\partial\psi}{\partial t}
+c\frac{\partial\psi}{\partial x}=0.
\label{eq:l9-transport-equation}
\end{equation}

\begin{figure}[htbp]
\centering
\includegraphics[width=0.88\textwidth]{lecture_09_figure_11.png}
\caption{The linear Hamiltonian reduces the Schr\"odinger equation to the first-order transport equation.}
\lectureframe{01:20:05}
\end{figure}

The general solution is
\begin{equation}
\psi(x,t)=f(x-ct),
\label{eq:l9-rigid-translation}
\end{equation}
for any normalizable initial profile \(f\). At \(t=0\), the profile is \(f(x)\). At later times, every feature has shifted to the right by \(ct\). The complex wave function, its real and imaginary parts, and its probability density all preserve their shape:
\begin{equation}
|\psi(x,t)|^2=|f(x-ct)|^2.
\end{equation}
Consequently,
\begin{equation}
\expect{X}_t=\expect{X}_0+ct.
\end{equation}

The symbol \(c\) was chosen to evoke the speed of light. The model resembles a single one-dimensional chiral massless excitation: every allowed component moves in one direction at one speed. It is not a complete photon theory. Low-energy one-dimensional phonons with a fixed sound speed provide another analogy. Because \(E=cp\), allowing both signs of \(p\) also allows negative energies, another sign that this is a deliberately restricted teaching model.

\begin{classroomqa}
\audiencequestion{Would \(H=cP\) describe any free particle if \(c\) were replaced by that particle's velocity?}
\lecturerresponse{No. A fixed coefficient in \(H=cP\) says that every momentum component has the same velocity. For a nonrelativistic particle the velocity is \(p/m\), so it depends on momentum. Writing \(vP\) with \(v=p/m\) is no longer linear in \(P\); it points toward a quadratic Hamiltonian.}
\end{classroomqa}

\begin{classroomqa}
\audiencequestion{Should the wave function of an ordinary freely moving electron not simply translate rigidly?}
\lecturerresponse{Only a very special dispersion relation produces rigid translation of every profile. An ordinary free-electron packet contains many momentum components whose phases evolve at different rates. Its center moves, but its shape generally spreads.}
\end{classroomqa}

\subsection{Classical comparison}

For the classical Hamiltonian \(H=cp\), Hamilton's equations give
\begin{equation}
\dot{x}=\frac{\partial H}{\partial p}=c,
\qquad
\dot{p}=-\frac{\partial H}{\partial x}=0.
\label{eq:l9-linear-hamilton-equations}
\end{equation}

\begin{figure}[htbp]
\centering
\includegraphics[width=0.88\textwidth]{lecture_09_figure_12.png}
\caption{Hamilton's equations for \(H=cp\): constant rightward velocity and conserved momentum.}
\lectureframe{01:25:20}
\end{figure}

The quantum probability distribution and its expectation value obey the same elementary motion. This agreement is the first concrete reminder that quantum evolution built from a classical Hamiltonian should reproduce the classical equations in the appropriate sense.

\section{The Nonrelativistic Free Particle}

For an ordinary particle of mass \(m\) with no force, classical kinetic energy is
\begin{equation}
H=\frac{p^2}{2m}.
\end{equation}
The quadratic form treats \(p\) and \(-p\) symmetrically, so left-moving and right-moving particles have the same positive energy. Quantization replaces \(p\) by the operator \(P\):
\begin{equation}
H=\frac{P^2}{2m}.
\label{eq:l9-free-hamiltonian}
\end{equation}
Substitute \(P=-i\hbar\,\partial_x\) into the Schr\"odinger equation:
\begin{align}
i\hbar\frac{\partial\psi}{\partial t}
&=
\frac{1}{2m}
\left(-i\hbar\frac{\partial}{\partial x}\right)^2
\psi \notag\\
&=
-\frac{\hbar^2}{2m}
\frac{\partial^2\psi}{\partial x^2}.
\label{eq:l9-free-schrodinger}
\end{align}

\begin{figure}[htbp]
\centering
\includegraphics[width=0.88\textwidth]{lecture_09_figure_13.png}
\caption{The free-particle Schr\"odinger equation obtained from \(H=P^2/(2m)\) and \(P=-i\hbar\,\partial_x\).}
\lectureframe{01:33:50}
\end{figure}

A plane-wave component with momentum \(p\) has energy
\begin{equation}
E(p)=\frac{p^2}{2m}.
\end{equation}
Unlike the linear case \(E=cp\), the slope of this dispersion relation depends on \(p\):
\begin{equation}
v_{\mathrm{group}}
=
\frac{dE}{dp}
=
\frac{p}{m}.
\end{equation}
Different momentum components of a localized packet therefore travel with different group velocities. Their relative phases change, the components separate, and the packet spreads. The decisive contrast is that \(H=cP\) transports every shape rigidly, whereas \(H=P^2/(2m)\) produces dispersive free-particle motion.

\section{What Has Been Established}

The mathematical interlude was not separate from the physics. Plane waves and the Dirac delta built the continuous basis machinery. The position operator acts by multiplication, the momentum operator acts by differentiation, and momentum eigenstates are plane waves. Projecting an arbitrary state onto those eigenstates produces the Fourier transform of its position wave function. Position and momentum are therefore rival representations of one state, and their uncertainty is the familiar reciprocal width of a Fourier pair.

The same representation machinery then turns a Hamiltonian into a differential equation. A linear function of momentum produces rigid translation; the quadratic kinetic energy produces the free-particle Schr\"odinger equation and dispersion. The course has moved from the foundations of measurement to the mechanics of motion without changing its central language: states, bases, Hermitian observables, amplitudes, and unitary evolution.
